\documentclass[10pt,a4paper]{amsart}
\usepackage[margin=19mm]{geometry}
\usepackage{amsmath,amssymb,mathtools}
\usepackage[scr=rsfs,cal=euler]{mathalfa}
\usepackage{xfrac}
\usepackage[unicode,hidelinks,bookmarks=false]{hyperref}
\newcommand{\ie}{i.e.\ }
\newcommand{\cf}{cf.\ }
\newcommand{\resp}{respectively\ }
\newcommand{\Subsection}{\subsection}
\providecommand{\nobreakdash}{\nobreak\hbox{-}}
\setlength{\emergencystretch}{2em}
\multlinegap0pt
\usepackage{amssymb}
\usepackage{mathtools}
\usepackage[shortlabels]{enumitem}
\newtheorem{lemma}{Lemma}
\newcommand{\C}{\mathbb{C}}
\newcommand{\eps}{\varepsilon}
\title{Optimal Linear Dependence on Boundary Type for Local Gromov Hyperbolicity of the Kobayashi Metric}
\author{Cheng Lou, Jianyong Qiao, Hongyu Wang, yumin Zhong}
\date{Source: arXiv:2609.15336v1 (CC BY 4.0)}
\begin{document}
\maketitle

\noindent\emph{Source and licence.} Optimal Linear Dependence on Boundary Type for Local Gromov Hyperbolicity of the Kobayashi Metric. Version arXiv:2609.15336v1 (CC BY 4.0), distributed by arXiv under the Creative Commons Attribution 4.0 International licence, http://creativecommons.org/licenses/by/4.0/, as recorded in the arXiv licence metadata for this exact version and shown on the abstract page https://arxiv.org/abs/2609.15336v1. Full licence notice: \texttt{licenses/arxiv-2609.15336-CC-BY-4.0.txt}.\par\smallskip

\noindent\textit{Editorial scope.} Counted unit: the article's lemma lem:extreme-infinite-type-point (source0.tex lines 5017--5023). The article's complete original proof of it is delivered with it (source0.tex lines 5025--5147, one \texttt{proof} environment of 123 lines). No mathematics is written, repaired or re-derived: the statement and the proof are the article's own lines, in the article's own notation, and no retained line is substituted or re-flowed.  No further source-local context is needed: every cross-reference of every retained line resolves inside the two delivered spans.  Nothing was omitted from any retained span, so the package carries no omission record.  No retained line contains a citation, so the wrapper carries no bibliography and no bibliography entry.  No retained line contains a figure environment, an image-inclusion command or drawing code, so no image is delivered and the package declares no asset dependency.  Editorial formatting only, in addition: an AMS \texttt{amsart} preamble that restates the article's own declarations and macros used by the retained lines, namely the environments \texttt{lemma} on separate counters, together with the standard packages those lines need.  The retained environments print in this document as Lemma 1 in order of appearance, whereas the article prints the same items with the numbering of its own full document.  The complete native source of the article as distributed in the arXiv e-print for this exact version is bundled unchanged as \texttt{sources/210410/source0.tex}.
\begin{lemma}
	\label{lem:extreme-infinite-type-point}
	Let \(\Omega\Subset\C^n\) be a convex domain with \(C^\infty\)-smooth
	boundary.  If \(\partial\Omega\) is not of finite line type, then there
	exists \(q\in\partial\Omega\) of infinite line type which is a complex
	extreme point of \(\overline\Omega\).
\end{lemma}

\begin{proof}
Let \(\rho\) be the signed distance function and consider the unit complex tangent bundle
\[
\mathcal S
=
\{(x,v):x\in\partial\Omega,\ v\in T_x^{\C}\partial\Omega,\ |v|=1\}.
\]
This bundle is compact.  For \(N\ge1\), let
\(\mathcal E_N\subset\mathcal S\) consist of the pairs \((x,v)\) satisfying
\[
\frac{\partial^{a+b}}
{\partial\zeta^a\partial\bar\zeta^b}
\rho(x+\zeta v)\bigg|_{\zeta=0}=0
\qquad
\text{whenever }1\le a+b\le N.
\]
Each \(\mathcal E_N\) is closed, and
\(\mathcal E_{N+1}\subset\mathcal E_N\).  Since the boundary has no
uniform finite line-type bound, every \(\mathcal E_N\) is nonempty. The Cantor intersection theorem therefore gives
\(\bigcap_{N\ge1}\mathcal E_N\ne\varnothing\).  Choose
\((p,v_0)\) in this intersection.  Then \(p\) has infinite line type.

If \(p\) is a complex extreme point of \(\overline\Omega\), there is
nothing to prove.  Otherwise, there exist
\(v\in\C^n\setminus\{0\}\) and \(\eps>0\)
	such that
	\[
	D:=\{p+\zeta v:|\zeta|<\eps\}\subset\overline\Omega.
	\]
	Let \(\ell\) be a real affine functional supporting
	\(\overline\Omega\) at \(p\), normalized so that
	\[
	\ell\le0\quad\text{on }\overline\Omega,
	\qquad
	\ell(p)=0.
	\]
	The map
	\(\zeta\longmapsto \ell(p+\zeta v)\)
	is real affine, nonpositive near \(0\), and vanishes at \(0\).  Hence it
	vanishes identically, and therefore
	\[
	D\subset
	F:=\overline\Omega\cap\{\ell=0\}
	\subset\partial\Omega.
	\]
	
	The set \(F\) is nonempty, compact, and convex.  Fix \(a\in\C^n\), and
	choose \(q\in F\) at which the function \(x\mapsto|x-a|^2\) attains its
	maximum on \(F\).  We claim that \(q\) is a real extreme point of \(F\).
	Indeed, if
	\[
	q=tx+(1-t)y,
	\qquad
	x,y\in F,\quad 0<t<1,
	\]
	with \(x\ne y\), then the strict convexity of \(|\,\cdot-a|^2\) gives
	\[
	|q-a|^2
	<
	t|x-a|^2+(1-t)|y-a|^2
	\le
	|q-a|^2,
	\]
	a contradiction.
	
	We next show that \(q\) is a real extreme point of
	\(\overline\Omega\).  Suppose that
	\[
	q=tx+(1-t)y,
	\qquad
	x,y\in\overline\Omega,\quad 0<t<1.
	\]
	Since
	\[
	0=\ell(q)=t\ell(x)+(1-t)\ell(y)
	\]
	and \(\ell(x),\ell(y)\le0\), we have
	\(\ell(x)=\ell(y)=0\).  Hence \(x,y\in F\), and the real extremality of
	\(q\) in \(F\) gives \(x=y=q\).  Thus \(q\) is a real extreme point of
	\(\overline\Omega\), and therefore also a complex extreme point of
	\(\overline\Omega\).
	
	It remains to show that \(q\) has infinite line type.  For \(0<t<1\),
	set
	\(p_t=(1-t)q+tp\).
	Since \(F\) is convex and contains both \(q\) and \(D\), it follows that
	\[
	\{p_t+\zeta v:|\zeta|<t\eps\}
	=
	(1-t)q+tD
	\subset F
	\subset\partial\Omega.
	\]
	
	Let \(\rho\) be a smooth defining function for \(\Omega\) in a
	neighborhood of \(q\).  Since \(p_t\to q\) as \(t\to0^+\), for all
	sufficiently small \(t>0\), the following identity holds:
	\[
	\rho(p_t+\zeta v)=0
	\qquad
	\text{for }|\zeta|<t\eps.
	\]
	Consequently, we have
	\[
	\frac{\partial^{a+b}}
	{\partial\zeta^a\partial\bar\zeta^b}
	\rho(p_t+\zeta v)\bigg|_{\zeta=0}
	=0
	\qquad
	\text{for all }a,b\ge0.
	\]
	Letting \(t\to0^+\) and using the smoothness of \(\rho\), we obtain
	\[
	\frac{\partial^{a+b}}
	{\partial\zeta^a\partial\bar\zeta^b}
	\rho(q+\zeta v)\bigg|_{\zeta=0}
	=0
	\qquad
	\text{for all }a,b\ge0.
	\]
	Thus \(\rho(q+\zeta v)\) vanishes to infinite order at \(0\), so \(q\)
	has infinite line type.
\end{proof}

\end{document}
